\section{¿Por qué regresan los modelos recurrentes/lineales}

En esta clase, aunque superficialmente se discute la compensación entre eficiencia de State Space Model (SSM, modelo de espacio de estados) y Transformer, la pregunta fundamental que realmente hay que responder es: \textbf{¿en qué forma guarda el pasado un modelo autoregresivo?} Si el pasado se almacena como una caché direccionable por token, el modelo excela en búsquedas precisas; si el pasado se comprime en un estado fijo, el modelo resulta más adecuado para actualizaciones en línea y abstracción. Las diferencias de complejidad entre ambos son simplemente el resultado externo de esa interfaz de memoria, no toda su esencia.

\subsection{De las corrientes de investigación a los sistemas de producción}

Primero, veamos por qué el ponente llama a esto «renacimiento». Tras Mamba surgieron múltiples rutas como xLSTM, DeltaNet, Gated DeltaNet y Test-Time Training; lo más importante es que ya han entrado en grandes modelos híbridos escalables: Jamba, Zamba, Samba, Hunyuan, Qwen, Kimi Linear, OLMo Hybrid y Megatron--NeMo. Al leer la figura, no hay que apresurarse a comparar nombres, sino identificar dos capas de información: a la izquierda están las familias de mecanismos; a la derecha, cómo entran esos mecanismos en modelos híbridos a escala de producción.

\lecturefigure{slide-002.jpg}{Familia de modelos recurrentes/lineales modernos y su adopción en grandes modelos híbridos}{Grabación oficial de Stanford, 00:03:20--00:04:09}

\teachervoice{00:01:28--00:03:55, el docente enfatiza que estos métodos ya han penetrado ampliamente en modelos a gran escala; por eso esta clase no presenta cada artículo individualmente, sino que asciende a características comunes, cuellos de botella compartidos y las diferencias fundamentales con Transformer.}

Un mismo mecanismo puede llamarse SSM, atención lineal, RNN moderno o modelo recurrente moderno según el comunidad. La siguiente figura coloca todos estos términos en una misma página; no significa que todas las fórmulas sean idénticas, sino que esta clase solo se preocupa por la interfaz compartida de \textbf{estado recurrente finito + cálculo secuencial lineal/cuadrático}. Por su parte, un modelo híbrido intercala capas de este tipo con capas de atención cuadrática en la dirección de profundidad de la red.

\lecturefigure{slide-003.jpg}{Taxonomía paraguas adoptada en esta clase: SSM, atención lineal y modelos recurrentes modernos}{Grabación oficial de Stanford, 00:04:10--00:05:09}

\begin{knowledgebox}{Términos clave: nombres distintos, ejes de comparación iguales}
\begin{tabularx}{\textwidth}{L{0.20\textwidth}X X}
\toprule
Término & Punto de partida habitual & Interfaz común de interés para esta clase \\
\midrule
SSM & Espacio de estados continuo/discreto, transición estructurada & Compresión de historia mediante recurrencia con estado finito \\
Atención lineal & Escribir el núcleo de atención como estadística recursiva & No guardar la matriz completa de atención $L\times L$ \\
RNN lineal & Actualización recurrente lineal y entrenamiento paralelo & Actualización de estado combinable y escaneable \\
Modelo recurrente moderno & Nueva generación de recurrencia respecto a LSTM/GRU & Ampliación del estado, fortalecimiento del gating y adaptación a aceleradores \\
Modelo híbrido & Intercalar diferentes tipos de capas desde la perspectiva de ingeniería & Capas recurrentes encargadas de la mayor parte del procesamiento, capas de atención para recuperación precisa \\
\bottomrule
\end{tabularx}
\end{knowledgebox}

\begin{warningbox}{No confundir la etiqueta paraguas con equivalencia matemática}
El formato de estado, las fórmulas de actualización, los algoritmos de entrenamiento y los kernels de Mamba, Gated DeltaNet, xLSTM y TTT son diferentes. Esta clase los agrupa en una sola categoría solo para comparar el eje principal entre «estado comprimido finito» y «caché a resolución de token»; cualquier conclusión sobre una implementación concreta debe volver a los artículos correspondientes y al hardware.
\end{warningbox}

\subsection{¿Por qué usar modelado autoregresivo como microscopio}

La generación autoregresiva descompone la distribución conjunta en probabilidades condicionales secuenciales: durante el entrenamiento se predice el siguiente token de forma paralela, pero en la inferencia hay que devolver el token recién generado al modelo. Por ello, lo que se deja «entre pasos» durante el proceso de decodificación expone el contrato de memoria real de la arquitectura.

\[
p(x_{1:L})=\prod_{t=1}^{L}p(x_t\mid x_{<t}),
\qquad
x_{t+1}\sim p_\theta(\cdot\mid x_{\le t}).
\]
Donde, $L$ es la longitud de la secuencia, $x_t$ es el $t$-ésimo token, $x_{<t}$ es el contexto previo y $p_\theta$ es el modelo con parámetros $\theta$. El entrenamiento puede procesar toda una secuencia a la vez; la generación debe avanzar en orden secuencial: $t=1,2,\ldots$.

\lecturefigure{slide-004.jpg}{Modelado autoregresivo: predicción paralela del siguiente token durante el entrenamiento, muestreo secuencial e inyección de retroalimentación durante la inferencia}{Grabación oficial de Stanford, 00:05:10--00:06:09}

\teachervoice{00:05:10--00:05:49, el ponente eligió el modelado autoregresivo por dos razones: es tanto el paradigma central de los modelos de lenguaje actuales como el que más fácilmente expone el estado y la complejidad de diferentes modelos secuenciales.}

\begin{importantbox}{La unidad de comparación de esta clase no es la capa, sino el estado de generación}
Al comparar dos arquitecturas, primero pregunte: después de completar el paso $t$ de generación, ¿qué debe conservar el modelo para generar el paso $t+1$? Ese objeto persistente entre pasos es el estado autoregresivo. Todas las discusiones posteriores sobre recuperación (recall), compresión, procesamiento en línea y resolución de token se desarrollan a partir de aquí.
\end{importantbox}

\subsection{Resumen de la sección}

El valor de los modelos recurrentes/lineales modernos ya no reside solo en reformular la complejidad en los artículos académicos, sino en el espacio de diseño real de los modelos de producción. Para atravesar la niebla de nombres de modelos, esta clase fija una lente de comparación: observar el estado que se preserva entre pasos durante la generación autoregresiva. En el capítulo siguiente se verá que Transformer guarda el pasado en una caché a resolución de token, mientras que el SSM lo comprime en un estado recurrente finito; ambas interfaces determinan simultáneamente las capacidades, los sesgos inductivos y los costos.
